\part{String-Net / State-Sum Models}

\section{Fusion categories}

\begin{exercise}[Closed-string configurations and a commuting-projector Hamiltonian]
Let $\Gamma$ be a connected trivalent graph giving a cell decomposition
of a closed oriented surface, with vertices $V$, edges $E$, and disc
faces $F$. A string configuration is a subset of $E$ with an even
number of occupied edges incident at every vertex. On
$\mathcal H=\C^{\{0,1\}^E}$ define
\[
 (Z_e\psi)(n)=(-1)^{n_e}\psi(n),\quad
 (X_e\psi)(n)=\psi(n+\mathbf1_e),\quad
 Q_v=\frac{1+\prod_{e\ni v}Z_e}{2},\quad
 B_f=\frac{1+\prod_{e\in\partial f}X_e}{2},
\]
where additions and incidences are modulo two. Assume each face
boundary is an embedded cycle. For $J_v,J_f>0$ put
$H=\sum_vJ_v(1-Q_v)+\sum_fJ_f(1-B_f)$.
\begin{enumerate}[label=(\alph*)]
\item Prove that the $Q_v$ impose the string constraint and that
$B_f$ averages a configuration with the configuration obtained by
adding the boundary of $f$. Prove that all these orthogonal
projections commute.
\item Prove that the zero-energy states are the functions supported
on string configurations and constant on orbits of face-boundary
addition. Construct one such state per orbit by uniform summation.
Prove invariance under every automorphism of the cell decomposition
that preserves the couplings.
\end{enumerate}
\end{exercise}

\begin{exercise}[Fusion categories and group-labelled strings]
A fusion category over $\C$ is a $\C$-linear semisimple rigid monoidal
category with finite-dimensional morphism spaces, finitely many simple
isomorphism classes, and simple tensor unit; semisimplicity and rigidity
have the meanings of the section Modular tensor categories. A spherical
structure is a pivotal structure whose left and right traces agree.
For a finite group $G$ let $\Vect_G$ consist of finite-dimensional
$G$-graded spaces, grading-preserving maps, and
$(U\otimes V)_g=\bigoplus_{ab=g}U_a\otimes V_b$.
\begin{enumerate}[label=(\alph*)]
\item Prove that $\Vect_G$ is a fusion category with simple objects
the specified graded lines $\C_g$, dual $\C_g^*=\C_{g^{-1}}$,
and fusion $\C_g\otimes\C_h=\C_{gh}$. Use the vector-space evaluation
and coevaluation to construct a spherical structure with all simple
dimensions equal to one.
\item For $G=\Z/2$ prove that the admissible trivalent labellings
$a+b+c=0$ are exactly the string configurations of the preceding
exercise. Interpret adding a face boundary as inserting the
nontrivial group-labelled loop and applying the fusion rule.
\end{enumerate}
\end{exercise}

\section{\texorpdfstring{$6j$}{6j}-symbols}

\begin{exercise}[Associators as maps between fusion spaces]
For simple objects $a,b,c,d$ in a fusion category define
$V_{ab}^c=\Hom(c,a\otimes b)$. Composition and the associator induce
the recoupling isomorphism
\[
 F^{abc}_d:
 \bigoplus_e V_{ab}^e\otimes V_{ec}^d
 \longrightarrow
 \bigoplus_f V_{bc}^f\otimes V_{af}^d.
\]
These intrinsic recoupling maps are the $6j$ data used here.
\begin{enumerate}[label=(\alph*)]
\item Construct this map by identifying the left and right sides
with $\Hom(d,(a\otimes b)\otimes c)$ and
$\Hom(d,a\otimes(b\otimes c))$, respectively. Prove both identifications
using the simple summands and composition pairings.
\item In $\Vect_G$, prove that the only nonzero channel has
$e=ab$, $f=bc$, $d=abc$ and that $F$ is the identity on the
canonical graded lines. Reverse the orientation of a string and
verify that its label changes to its group inverse.
\item Replace the associator on the canonical lines by multiplication
by $\omega(a,b,c)\in U(1)$. Compute the inverse recoupling map
and show that it is the adjoint. Determine the normalization of
$\omega$ forced by keeping the unit constraints equal to the identity.
\end{enumerate}
\end{exercise}

\section{Pentagon identities}

\begin{exercise}[The pentagon, three-cocycles, and changes of associator]
A normalized group three-cocycle is a function $\omega:G^3\to U(1)$,
equal to one if an argument is the identity, satisfying
\[
 \omega(b,c,d)\omega(a,bc,d)\omega(a,b,c)
 =\omega(ab,c,d)\omega(a,b,cd).
\]
For a normalized function $\nu:G^2\to U(1)$ put
$ (\delta\nu)(a,b,c)=
\nu(b,c)\nu(a,bc)/(\nu(ab,c)\nu(a,b))$.
\begin{enumerate}[label=(\alph*)]
\item Compare the two composites from
$((a\otimes b)\otimes c)\otimes d$ to
$a\otimes(b\otimes(c\otimes d))$.
Prove that the associator of the preceding exercise satisfies
the pentagon exactly when $\omega$ is a three-cocycle.
\item Prove that replacing $\omega$ by $\omega\,\delta\nu$
gives a monoidally equivalent category by giving the tensor
constraint of the identity functor. Denote this category by
$\Vect_G^\omega$.
Construct its left and right duals by multiplying the
evaluation and coevaluation on the canonical graded lines
by unit scalars, and use the cocycle equation to check
both snake identities.
\item On $G=\Z/2$ verify that $\omega(a,b,c)=(-1)^{abc}$
is a normalized three-cocycle and cannot equal $\delta\nu$ for
a normalized $\nu$, by evaluating at $(1,1,1)$.
\end{enumerate}
\end{exercise}

\section{Turaev--Viro}

\begin{exercise}[The pointed Turaev--Viro sum and local retriangulations]
Let $M$ be a closed oriented triangulated three-manifold, with an
ordering of its vertices. Label an oriented edge $ij$ by $g_{ij}\in G$,
with $g_{ji}=g_{ij}^{-1}$, and impose $g_{ij}g_{jk}=g_{ik}$ on each
ordered triangle. In the untwisted category $\Vect_G$ define
\[
 Z_G(M)=|G|^{-n_0}
 \sum_{\text{admissible edge labellings}}1,
\]
where $n_0$ is the number of vertices. This is the pointed,
untwisted Turaev--Viro state sum in this exercise.
\begin{enumerate}[label=(\alph*)]
\item For a $2$--$3$ move, compare two and three tetrahedra filling
the same triangular bipyramid. Prove that the triangle constraints
give a bijection between internal labellings for each admissible
boundary labelling.
\item For a $1$--$4$ move, inserting one vertex in a tetrahedron,
prove that an admissible boundary labelling has exactly $|G|$
extensions, parametrized by the label of one new edge.
Deduce invariance of $Z_G$ under both moves and their inverses.
\item Insert the oriented recoupling weight
$\omega(g_{ij},g_{jk},g_{k\ell})^{\epsilon_{ijk\ell}}$
on a tetrahedron $i<j<k<\ell$. Verify that equality of the products
in the $2$--$3$ move is precisely the pentagon identity.
Keep the subsequent state-space calculation untwisted.
\end{enumerate}
\end{exercise}

\section{TQFT}

\begin{exercise}[Finite field groupoids and gluing of state spaces]
A groupoid is a category in which every morphism is invertible.
For a finite groupoid $\mathcal G$ define
$|\mathcal G|=\sum_{[x]}|\Aut(x)|^{-1}$. For a triangulated
surface $\Sigma$, let $\mathcal A_\Sigma$ be the groupoid of
the admissible $G$-edge labellings of the preceding section;
a morphism is a vertex labelling $(h_i)$ acting by
$g_{ij}\mapsto h_i g_{ij}h_j^{-1}$. Set
$\mathcal V(\Sigma)=\Fun(\pi_0\mathcal A_\Sigma,\C)$ with inner product
$\sum_{[a]}\overline{f(a)}g(a)/|\Aut(a)|$.
An oriented bordism $M:\Sigma_0\to\Sigma_1$ is a compact
oriented three-manifold with boundary identified with
$\overline{\Sigma_0}\sqcup\Sigma_1$.
A three-dimensional oriented TQFT assigns spaces to closed oriented
surfaces and linear maps to oriented bordisms, compatibly with
composition and disjoint union.
\begin{enumerate}[label=(\alph*)]
\item Prove the orbit-counting identity
$|X//G|=|X|/|G|$ for a finite group action.
Recover $Z_G(M)$ as the cardinality of the field groupoid on $M$.
\item For a bordism $M:\Sigma_0\to\Sigma_1$ define its map by
pullback along restriction to $\Sigma_0$ followed by summation over
the homotopy fibres of restriction to $\Sigma_1$, weighted by
$|\Aut|^{-1}$. A homotopy fibre over $b$ consists of a field
and an isomorphism from its restriction to $b$.
Show by finite double counting that composition corresponds to
gluing, the cylinder gives the identity, and disjoint union gives
the tensor product. Work with fixed triangulated bordisms and the
local moves proved in the section Turaev--Viro.
\end{enumerate}
\end{exercise}

\begin{exercise}[Prediction: string-net ground-state degeneracy and excitation energy]
Return to the $\Z/2$ string Hamiltonian of the section Fusion
categories on a genus-$g$ surface, with constant $J_v,J_f>0$.
\begin{enumerate}[label=(\alph*)]
\item Prove that the even edge subsets form a vector space over
$\mathbb F_2$ of dimension $|E|-|V|+1$ by eliminating edges of
a spanning tree. Prove that face boundaries span a space of
dimension $|F|-1$: in a sum with zero boundary, coefficients
of adjacent faces agree. Using $|V|-|E|+|F|=2-2g$, derive
\[
 \dim\ker H=2^{2g}.
\]
\item The global products of the vertex involutions and of the
face involutions are both one. Deduce that each type of
violation occurs in pairs, and construct a string operator
creating a pair of either type. On a sufficiently large cell
decomposition derive the excitation gap $2\min(J_v,J_f)$.
\item Identify the degeneracy with $\dim\mathcal V(\Sigma)$ for
$G=\Z/2$ by exchanging the primal and dual cell decompositions.
State the predicted values on the sphere and torus and the
spectroscopic pair-creation thresholds $2J_v$, $2J_f$.
Specify that these are exact predictions for the stated
commuting-projector Hamiltonian; a microscopic material requires
an identification of its effective couplings and degrees of freedom.
Compare the construction with Levin and Wen \cite{levin-wen}.
\end{enumerate}
\end{exercise}
